\section{使命与可信 AI：从聪明到可信}

本章收束到吴明辉的个人目标。他说，GPT 出来之前，自己的人生目标是做一个比自己更聪明的 AI；后来他发现“聪明”已经被更大的模型公司推进了，于是目标转向做最可信的 AI，即 trusted agentic model。这个转向很关键：企业级 AI 的核心不是谁最会聊天，而是谁能被信任地用于真实业务。

\begin{figure}[H]
\centering
\includegraphics[width=0.96\textwidth]{happy-founder-mission.png}
\caption{幸福 Founder 的使命对齐：个人使命、家庭使命和公司使命高度相关。自制概念图，依据 02:53:20--03:47:40 对谈内容整理。}
\end{figure}

\begin{knowledgebox}{读图：使命对齐是长期投入的能源}
个人使命决定创始人愿意解什么题，家庭使命提供稳定能量，公司使命连接客户价值。吴明辉把创业困难看成数学题，背后是强延迟满足和使命对齐。
\end{knowledgebox}

\subsection{从最聪明到最可信}

前面讲组织和生产关系，本节回到企业 AI 的最终评价标准。ToC 模型竞争常常围绕“谁更聪明”；ToB 企业服务更关心“谁更可信”。可信包括数据来源可信、权限可信、过程可审计、结果可复核、错误可追责、系统可私有化或受控部署。企业客户不会只为炫技买单，最终要为风险可控的结果付费。

\begin{importantbox}{Trusted Agentic Model 的含义}
Trusted Agentic Model 是能在企业私有上下文中执行任务、调用工具、保留审计轨迹、尊重权限边界并输出可复核结果的智能体模型。它的目标不是单次回答最惊艳，而是长期运行最可信。
\end{importantbox}

\begin{knowledgebox}{可信 Agent 的验收清单}
\begin{tabularx}{\textwidth}{p{0.20\textwidth}p{0.32\textwidth}X}
\toprule
维度 & 检查问题 & 为什么重要 \\
\midrule
权限 & 它是否只读取和修改允许范围内的数据？ & 防止越权和合规事故。 \\
审计 & 每一步动作是否有日志和可复盘证据？ & 企业需要追责和改进。 \\
复核 & 高风险动作是否有人类确认？ & 降低不可逆错误。 \\
反馈 & 结果是否回流到模型和流程？ & 让系统持续变好。 \\
\bottomrule
\end{tabularx}
\end{knowledgebox}

\subsection{幸福的老板：把 suffering 当题解}

上一节讲可信 AI 的产品目标，本节回到支撑这种长期目标的人。结尾处吴明辉说自己是幸福的 founder，即便 2022 年很难，也没有特别 suffering，因为相信自己一定能搞定，只是解题时间更长。这不是轻松，而是一种把长期困难抽象为题目的心态。它解释了为什么 ToB 创业能熬 19 年：没有这种延迟满足和心理结构，很难穿过漫长交付、资本波动和组织重构。

\begin{knowledgebox}{ToB Founder 的心理结构}
ToB 创业很少有短期爆发的反馈，更多是长期客户关系、现金流约束、组织稳定和产品信用。能否把困难看成可分解问题，而不是身份失败，决定了 founder 是否能持续迭代。
\end{knowledgebox}

这也解释了为什么访谈最后不只是商业复盘，而是使命复盘。个人使命、家庭使命和公司使命高度相关，意味着创始人不是把公司当成短期项目，而是把公司当成自己长期问题的一部分。对企业级 AI 公司尤其如此：客户信任、数据安全、行业交付和可信 Agent 都需要长周期积累，无法只靠一轮模型能力窗口完成。

\begin{importantbox}{长期主义的真实成本}
长期主义不是口号，而是愿意承受慢反馈、低外部声量、现金流压力和组织反复调整。ToB 公司如果最终能进入 AI 时代，靠的往往不是某次爆发，而是十几年积累下来的客户、数据、交付和 founder 心理结构。
\end{importantbox}

\subsection{本章小结}

吴明辉的 AI 目标从“更聪明”转向“更可信”，恰好对应企业级 AI 的核心需求。企业客户需要的不只是模型能力，而是可信执行和长期交付。

